% !TeX program = lualatex
% Definition and sources only; this file is not an LLM solution attempt.
\documentclass[11pt,a4paper]{article}
\usepackage[margin=25mm]{geometry}
\usepackage{fontspec}
\setmainfont{lmroman10-regular.otf}[BoldFont=lmroman10-bold.otf,
 ItalicFont=lmroman10-italic.otf,BoldItalicFont=lmroman10-bolditalic.otf]
\usepackage{amsmath,amssymb,mathtools}
\usepackage{unicode-math}
\setmathfont{latinmodern-math.otf}
\AtBeginDocument{\let\setminus\smallsetminus}
\usepackage{xurl,hyperref}
\hypersetup{colorlinks=true,urlcolor=blue}
\setcounter{secnumdepth}{0}
\setlength{\parindent}{0pt}
\setlength{\parskip}{5pt}
\setlength{\emergencystretch}{3em}
\begin{document}
\section*{Some Questions — Question 3}\label{statement}

\subsection{Definitions and mathematical statement}
Let $T$ be a rooted locally finite tree without terminal vertices and with infinitely many boundary rays. Its boundary $\partial T$ consists of infinite rays starting at the root, with the compact topology in which rays are close if they share a long initial segment. Give $\operatorname{Aut}(T)$ the topology of agreement on finite rooted subtrees. Let $G\leq\operatorname{Aut}(T)$ be a closed subgroup whose action on $\partial T$ is transitive.

For an integer $m\geq0$, the level stabilizer is the subgroup fixing every vertex at distance $m$ from the root:
\[
 \operatorname{St}_G(m)=\{g\in G:gv=v\text{ for every level-}m\text{ vertex }v\}.
\]
Must this subgroup contain an element $g$ that fixes no boundary ray? In quantified form the question is whether
\[
 \forall m\geq0\ \exists g\in\operatorname{St}_G(m)\quad
                  \operatorname{Fix}_{\partial T}(g)=\varnothing.
\]
Fixed-point-freeness refers to the boundary, not to vertices of the tree: every automorphism fixes the root, and a level stabilizer necessarily fixes the first $m$ levels. The same selected element must move every infinite ray, rather than choosing a different element for each ray. Closedness is part of the hypothesis, making $G$ a compact group. Transitivity likewise means actual transitivity on complete rays, not merely that an individual orbit is dense.
\subsection{Short English statement}
Must every level stabilizer contain a single automorphism that moves every boundary ray?
\subsection{Sources}
\begin{itemize}
\item[S1] ULAM.AI and the UnsolvedMath Contributors, supplied problems.json, catalogue entry, AMR Open Problem Lists. Catalogue identity and source transcription; CC BY 4.0.
\url{https://huggingface.co/datasets/ulamai/UnsolvedMath}.
\item[S2] Abert - Some questions (pdf) (2010), Question 3. 
\url{https://www.renyi.hu/~abert/questions.pdf}.
\end{itemize}
\end{document}
